% § 9 — Discussion and Limitations
\section{Discussion and Limitations}
\label{sec:discussion}
Three points follow. First, the grey list is executable, and executing it literally is precise but narrow. On five of the nine scored items, at least six flags in ten are correct, but the configurations the Annex describes cover a small part of what the \corpus{} annotators labelled: our queries retrieve a third of it, and a fine-tuned encoder more than twice as much. Second, the clause theme carries more retrievable signal than the templates add. What the templates contribute is \emph{grounds}---the ability to say which item fired and why---rather than signal: the only significant gain, on one item, does not survive a second extraction pass. The error analysis shows where the price of these grounds is paid: fields that block a query and actions missing from the templates, not absent templates. Third, the mismatch between the Annex and the benchmark is itself a finding. Categories such as choice of law have no grey-list ground; items such as \gi{d}, \gi{m} and \gi{p} fire on the held-out contracts but have no \corpus{} counterpart, so a detector that executed the law faithfully would be penalised by the benchmark whenever it is right. \textbf{Limitations and future work.} The templates were not reviewed by legal experts within the time frame of this study. As a consequence, low recall on a given item cannot be attributed with certainty to the query's narrowness rather than to extraction noise on the relevant field; the two are conflated in this end-to-end evaluation and will be disentangled by the expert validation described below. Our results characterise the pipeline as a whole and do not establish the correctness of individual templates. The choice of extractor matters: on the development pilot, with the previous prompt version, a second model (gpt-6-astra) agreed with the one used here on the decisive fields with a mean Jaccard index of about 0.6. The held-out stability estimate rests on a single pair of extraction runs. The reference labels are sentence-level while the Annex is clause-level; thresholds for quantitative items are declared, not derived; the held-out set is small and the intervals are wide; the queries reflect one reading of one legal system's list. This work is part of an ongoing collaboration with legal experts: the legal co-authors helped design the architecture of the system and brought the legal practitioner's perspective to it. The follow-up study addresses these limitations: independent double validation of the templates by two legal experts, with per-field agreement, which will separate extraction errors from query errors; re-extraction with pinned model versions and replication with an open-weight model; and an expert audit of the flagged sentences and of the explanations the queries return.
